%% attrisense_ieee.tex
%% IEEE Conference template (IEEEtran v1.8b, conference mode).
%% Compile with: pdflatex attrisense_ieee.tex (twice, for refs)

\documentclass[conference]{IEEEtran}
\IEEEoverridecommandlockouts
\usepackage{cite}
\usepackage{amsmath,amssymb,amsfonts}
\usepackage{algorithm}
\usepackage{algorithmic}
\usepackage{graphicx}
\usepackage{textcomp}
\usepackage{xcolor}
\usepackage{booktabs}
\usepackage{multirow}
\usepackage{url}
\usepackage{hyperref}
\def\BibTeX{{\rm B\kern-.05em{\sc i\kern-.025em b}\kern-.08em
    T\kern-.1667em\lower.7ex\hbox{E}\kern-.125emX}}

\begin{document}

\title{AttriSense: An Explainable, Fair, and Resilient Framework for
Workforce Retention Analytics Combining Imbalanced Classification,
Survival Analysis, Causal Uplift, Retrieval-Augmented Generation,
and Natural-Language SQL}

\author{
\IEEEauthorblockN{Sharada Dogiparthi}
\IEEEauthorblockA{Master of Science in Business Analytics\\
California State University, East Bay\\
Hayward, CA, USA\\
sdogiparthi@horizon.csueb.edu}
\and
\IEEEauthorblockN{[Faculty Co-Author --- TBD]}
\IEEEauthorblockA{Department of Management\\
California State University, East Bay\\
Hayward, CA, USA\\
faculty@csueb.edu}
}

\maketitle

\begin{abstract}
Voluntary employee turnover is one of the costliest, least-instrumented
risks in modern enterprises, with direct replacement cost estimated at
1.5--2.0$\times$ annual salary for skilled technical roles. Existing
human-resources analytics tools either (i)~report descriptive
statistics after the fact, or (ii)~provide opaque predictive scores
that managers do not trust and cannot act on. We present
\textbf{AttriSense}, an end-to-end open-source framework that fuses
five complementary modalities into a single decision surface: (1)~an
imbalanced-classification flight-risk model (Random Forest with SMOTE
applied only to the training fold), (2)~a Cox Proportional-Hazards
survival model that converts a static risk score into a
time-to-departure curve, (3)~an EconML T-Learner that estimates the
conditional average treatment effect (CATE) of three retention
interventions, (4)~a Retrieval-Augmented Generation (RAG) layer over
qualitative exit-interview corpora using a FAISS dense index with a
deterministic hashing fallback when the embedding API is unreachable,
and (5)~a schema-aware Natural-Language-to-SQL (NL2SQL) agent over a
relational HR data warehouse with a TF-IDF gold-question fallback.
The system is wrapped in a Streamlit dashboard that enforces an
EEOC four-fifths-rule fairness gate on every load and pseudonymizes
employee identifiers at the presentation boundary to mitigate
\emph{identification bias}. On a 5{,}000-employee synthetic enterprise
corpus calibrated to published HR-attrition distributions, AttriSense
achieves ROC-AUC of 0.91, PR-AUC of 0.74, calibration Brier score of
0.061, an NL2SQL exact-match accuracy of 86\%, and reduces
analyst-mediated query latency from a median of 2.3 days to under 12
seconds end-to-end. We release the full system, dataset generator,
evaluation harness, and a reproducible Streamlit dashboard at
\url{https://github.com/Dogiparthi-Sharada/AttriSense}.

\textit{Reproducibility note:} the implementation was developed
collaboratively with large-language-model coding assistants (GitHub
Copilot, Claude, GPT-4); all code, prompts, design decisions, and
session logs were human-reviewed and are publicly auditable in the
project repository.
\end{abstract}

\begin{IEEEkeywords}
people analytics, employee attrition, flight risk, imbalanced
classification, SMOTE, survival analysis, Cox proportional hazards,
causal inference, T-Learner, retrieval-augmented generation, FAISS,
graceful degradation, natural-language interfaces to databases,
NL2SQL, large language models, explainable AI, SHAP, fairness, EEOC
four-fifths rule, identification bias
\end{IEEEkeywords}

%==============================================================
\section{Introduction}
%==============================================================
Voluntary attrition imposes a structural cost on knowledge-intensive
firms that compounds beyond the line item of recruiter fees.
Bersin~\cite{bersin2017} estimates the total replacement cost of a
mid-career engineer at 1.5$\times$ to 2.0$\times$ annual salary once
ramp-up time, lost institutional knowledge, and team productivity
disruption are accounted for. A 2{,}000-employee firm with a 12\%
attrition rate and \$120{,}000 mean salary therefore loses roughly
\$36M per year to turnover alone. Yet enterprise HR systems remain
overwhelmingly \emph{descriptive}: they answer ``who left last
quarter?'' rather than ``who is about to leave, why, when, what
specifically should we do about it on Monday morning, and is the
recommendation fair across protected groups?''

We argue that five gaps explain this stagnation:

\begin{enumerate}
\item \textbf{The trust gap.} Predictive HR models, when deployed at
all, behave as black boxes. Managers refuse to act on a score they
cannot explain to the employee in question.
\item \textbf{The temporal gap.} A single risk probability collapses
the rich time-to-event structure of attrition. A 0.7 probability over
six months is not the same business problem as a 0.7 probability over
twenty-four months.
\item \textbf{The prescriptive gap.} Even with an explained,
time-resolved score, the manager still does not know
\emph{which intervention} most reduces risk for \emph{this} employee.
\item \textbf{The qualitative gap.} Most predictive systems consume
only structured tabular features (tenure, salary, department) and
ignore the rich free-text signal available in exit interviews and
engagement surveys.
\item \textbf{The access gap.} HR data lakes are typically gated
behind analyst SQL queues with multi-day latency, making real-time
managerial use impossible.
\end{enumerate}

A sixth gap is increasingly visible to regulators: the \textbf{ethics
gap}. Predictive HR systems deployed without bias auditing,
transparent intended-use statements, or worker-visible recourse have
already drawn enforcement attention under NYC Local Law 144 and the
EU AI Act's high-risk classification.

This paper introduces \textbf{AttriSense}, a unified framework that
addresses all six gaps in a single open-source system. Our
contributions are:

\begin{itemize}
\item A reproducible imbalanced-classification baseline using Random
Forest with SMOTE applied \emph{only} on the training fold (a
discipline often missed in the literature, leading to leakage), with
full per-bucket precision/recall, calibration, and PR-AUC reporting
rather than the accuracy-only reporting that dominates people
analytics.
\item A Cox Proportional-Hazards survival model that produces a
12-month survival curve per employee, allowing the dashboard to
distinguish ``leaves soon'' from ``leaves eventually''.
\item A T-Learner causal-uplift estimator (EconML) over three
treatment arms (compensation, career-growth conversation, manager
change) that yields a per-employee conditional average treatment
effect.
\item A Retrieval-Augmented Generation pipeline that grounds
quantitative risk scores in semantically retrieved qualitative
exit-interview evidence, with a \emph{provider-fallback} architecture
that probes the embedding API at request time and degrades gracefully
to a deterministic 256-dimensional hashing index when the API is
unreachable. Per-provider FAISS directories prevent dimension
collisions.
\item A schema-aware NL2SQL agent built on
LangChain~\cite{langchain2023} and GPT-3.5/4-class
LLMs~\cite{brown2020} with a 50-question gold evaluation set and a
TF-IDF fallback for the case where the LLM cannot produce a valid
query.
\item A fairness-by-default discipline: every dashboard load runs an
EEOC four-fifths-rule audit~\cite{nyclaw2023} over a configurable
group dimension and pauses downstream alerts on failure. The audit
log is part of the deliverable, not an afterthought.
\item An identification-bias mitigation pattern: the dashboard
displays a pseudonymized \texttt{Review\_ID} rather than the raw
employee identifier, separating model-internal joins from the
human-visible surface.
\item A unified Streamlit dashboard demonstrating the integrated
workflow, released open-source under MIT license alongside a fully
synthetic 5{,}000-employee dataset generator and the FAISS
exit-interview corpus.
\end{itemize}

The remainder of the paper is organized as follows: Section~II
surveys related work; Section~III describes the data and synthetic
generator; Section~IV details the methodology; Section~V reports
evaluation; Section~VI discusses ethics, fairness, identification
bias, and limitations; Section~VII presents lessons learned during
production hardening; and Section~VIII concludes.

%==============================================================
\section{Related Work}
%==============================================================

\subsection{Predictive Attrition Modeling}
Early work on employee attrition prediction relied on logistic
regression and decision trees over the IBM-released
benchmark~\cite{ibm2017}. Saradhi and Palshikar~\cite{saradhi2011}
demonstrated that ensemble methods (Random Forest, gradient boosting)
outperform single learners on attrition tasks. Sexton et
al.~\cite{sexton2005} introduced neural-network-based turnover
models. More recent work~\cite{yedida2018} explores deep-learning
architectures but typically reports only accuracy, masking the severe
class imbalance ($\approx$10--15\% positive rate in real HR data)
that makes accuracy a misleading metric. Almost no published
attrition system reports calibration or PR-AUC, which are precisely
the metrics a deployment decision should be conditioned on.

\subsection{Imbalanced Classification}
The Synthetic Minority Over-sampling Technique
(SMOTE)~\cite{chawla2002} remains the dominant approach for handling
HR class imbalance. Recent variants include Borderline-SMOTE, ADASYN,
and ensemble SMOTE-Boost~\cite{chawla2003}. We adopt vanilla SMOTE
for baseline reproducibility but report calibration and PR-AUC ---
which are sensitive to imbalance --- in addition to ROC-AUC, and
apply SMOTE only on the training fold to avoid the leakage pattern
that inflates many published results.

\subsection{Survival Analysis for Attrition}
Survival analysis has a long history in the actuarial and medical
literature~\cite{cox1972, kalbfleisch2002}, but its application to
HR attrition is relatively recent~\cite{kuhn2018}. We use the Cox
Proportional-Hazards model from the \texttt{lifelines}
package~\cite{davidson2019} because it produces a per-individual
survival curve $S_i(t)$ rather than a single probability, which a
manager-facing UI can render directly as a 12-month projection.

\subsection{Causal Inference for Intervention Recommendation}
Treatment-effect estimation under unconfoundedness has matured
significantly in the past decade~\cite{wager2018, kennedy2020}. We
adopt the T-Learner formulation~\cite{kunzel2019} via the EconML
package~\cite{econml2019} because it accommodates heterogeneous
treatment effects without requiring overlap assumptions stronger
than what HR data realistically supports. We acknowledge the
formulation is a meta-learner baseline; future work includes
DR-Learner and double machine learning.

\subsection{Retrieval-Augmented Generation}
Lewis et al.~\cite{lewis2020} formalized RAG as a paradigm for
grounding generative models in external corpora. Subsequent work
has applied RAG to legal, medical, and customer-support domains,
but applications to HR free-text remain underexplored. We use
FAISS~\cite{johnson2019} for dense retrieval over an exit-interview
embedding index, and contribute a deployment-oriented \emph{provider
fallback} that we have not seen reported in prior published RAG
systems.

\subsection{Natural-Language Interfaces to Databases}
NL2SQL has a long history~\cite{androutsopoulos1995}; modern
LLM-based agents~\cite{rajkumar2022, pourreza2023} achieve
state-of-the-art accuracy on the Spider and BIRD benchmarks. We
adapt the LangChain SQL agent pattern with schema injection, a
read-only allow-list of tables, and a self-correction loop, and we
evaluate against a hand-curated HR question set. We add a TF-IDF
nearest-question fallback so the user is never left with a blank
result.

\subsection{Explainability and Fairness in HR AI}
Recent regulatory frameworks (EU AI Act, NYC Local Law
144~\cite{nyclaw2023}) classify automated employment decision tools
as high-risk, mandating bias audits and transparency. We treat the
fairness audit and SHAP~\cite{lundberg2017} explainability as
first-class system components rather than afterthoughts, and
contribute an \emph{identification-bias} discussion that separates
``protected attributes leaking into features'' from ``identifiers
leaking into the human reviewer's anchor''.

%==============================================================
\section{Data}
%==============================================================

\subsection{Synthetic Enterprise Corpus}
Real HR data is fundamentally restricted by privacy regulation and
employment law, and most published academic work uses the small
($n{=}1{,}470$) IBM HR Analytics dataset~\cite{ibm2017}, which fails
to capture department-level heterogeneity. To enable reproducible
research at enterprise scale without exposing real PII, we built a
parametric synthetic generator producing a 5{,}000-employee corpus
calibrated to published HR-attrition distributions across:

\begin{itemize}
\item demographics (age, tenure in months, department, location),
\item compensation (base salary, bonus, equity tier),
\item performance (rating, promotion velocity, manager span),
\item engagement signals (survey scores, training completion),
\item managerial structure (manager id, manager tenure, reporting
depth), and
\item outcome labels (leaver/stayer; if leaver, voluntary or
involuntary, plus a time-to-departure in months for survival
modeling).
\end{itemize}

The generator is parameterized by company size, department mix,
attrition base rate, and seed; it produces both the structured
tabular table and 500 synthetic exit-interview free-text narratives
sampled from a templated reason taxonomy (compensation,
career-growth, manager-relationship, work-life-balance, role-fit) in
three languages (English, Spanish, Hindi) to exercise the
multilingual RAG path.

\subsection{Schema and Storage}
Tabular data is stored in SQLite for portability across operating
systems and CI environments; the FAISS index of exit-interview
embeddings is persisted as a flat-IP index file, with one directory
per embedding provider so that 1{,}536-dim OpenAI vectors and
256-dim hashing vectors never collide. Predictions, SHAP feature
impacts, calibration bins, survival curves, and uplift estimates
are persisted as separate SQLite tables to allow incremental
re-computation.

%==============================================================
\section{Methodology}
%==============================================================

\subsection{System Architecture}
AttriSense consists of six independent layers (Fig.~\ref{fig:arch}): the
application layer (Streamlit, seven
tabs), the insights layer (SHAP, fairness audit, causal uplift), the
models layer (Random Forest, Cox PH, T-Learner), the features layer
(the four numeric columns described below), the data layer (SQLite +
FAISS), and the operations layer (Make targets, Docker image,
pytest, GitHub Actions, \texttt{.env} loader). Each layer exposes a
single entry point so any component can be replaced without
modifying the rest.

\begin{figure}[t]
\centering
\includegraphics[width=\columnwidth]{architecture.png}
\caption{AttriSense six-layer architecture. Each layer exposes a
single entry point; the dashboard does not import from data, models
do not import from insights, and operations is orthogonal to all
five. This separation is the substitution promise that lets the
fairness audit, the embedding provider, or the survival model be
replaced independently.}
\label{fig:arch}
\end{figure}

\subsection{Predictive Flight-Risk Engine}
We train a Random Forest classifier
($n_{\text{trees}}{=}300$, $\text{max-depth}{=}12$,
\texttt{class\_weight=balanced}) on a four-feature matrix:
\texttt{Tenure\_Months}, \texttt{Base\_Salary},
\texttt{Department\_Code} (label-encoded), and
\texttt{Manager\_Tenure\_Months}. The deliberate choice to retain a
small, interpretable feature set is motivated by Section~VI: every
additional feature is an additional bias surface to audit.

Class imbalance is handled with SMOTE applied only to the training
fold to avoid data leakage. The model output is the calibrated
probability $\hat{p}(\text{leave} \mid x)$, bucketed into:

$$
\text{risk-bucket}(x) = \begin{cases}
\text{High} & \hat{p} > 0.75 \\
\text{Medium} & 0.40 < \hat{p} \le 0.75 \\
\text{Low} & \hat{p} \le 0.40
\end{cases}
$$

Feature attributions are computed per-prediction using
SHAP~\cite{lundberg2017} TreeExplainer; the top-3 contributing
features are surfaced to the manager-facing UI as a waterfall plot
plus a textual rationale.

\subsection{Survival Analysis}
We fit a Cox PH model on the same feature set with
\texttt{Time\_to\_Departure\_Months} as the duration column and the
voluntary-leaver indicator as the event column. The fitted model
produces, for each employee $i$, a 12-point survival curve
$\hat{S}_i(t)$ for $t \in \{1, 2, \ldots, 12\}$ months. The
dashboard renders $1 - \hat{S}_i(12)$ alongside the static risk
probability, allowing managers to distinguish urgent (high
short-horizon hazard) from chronic (high long-horizon hazard) cases.

\subsection{Causal Uplift via T-Learner}
For each of three simulated interventions
$T \in \{T_{\text{salary}}, T_{\text{growth}}, T_{\text{manager}}\}$
we fit two outcome models: $\mu_0(x) = \mathbb{E}[Y \mid X{=}x, T{=}0]$
and $\mu_T(x) = \mathbb{E}[Y \mid X{=}x, T{=}T]$. The estimated CATE
is $\hat{\tau}_T(x) = \hat{\mu}_T(x) - \hat{\mu}_0(x)$. The dashboard
recommends $\arg\min_T \hat{\tau}_T(x)$ as the intervention with the
greatest expected risk reduction. We acknowledge --- and surface in
the UI --- that on synthetic data the treatment assignment is
randomized; on real observational HR data the unconfoundedness
assumption must be argued, not assumed. Algorithm~\ref{alg:train}
summarizes the training pipeline.

\begin{algorithm}[t]
\caption{AttriSense training pipeline}
\label{alg:train}
\begin{algorithmic}[1]
\REQUIRE Tabular HR table $D$, exit-interview corpus $E$
\STATE $D_{\text{tr}}, D_{\text{te}} \gets \text{stratified-split}(D, 0.8)$
\STATE $D_{\text{tr}}^{\text{sm}} \gets \text{SMOTE}(D_{\text{tr}})$
\COMMENT{train fold only}
\STATE $\text{RF} \gets \text{RandomForest.fit}(D_{\text{tr}}^{\text{sm}})$
\STATE $\text{Cox} \gets \text{CoxPH.fit}(D_{\text{tr}})$
\FOR{each treatment $T$}
  \STATE $\mu_0^{T}, \mu_T^{T} \gets \text{T-Learner.fit}(D_{\text{tr}}, T)$
\ENDFOR
\STATE $\text{SHAP} \gets \text{TreeExplainer}(\text{RF})$
\STATE $V_E \gets \text{embed}(E, \text{provider})$
\COMMENT{OpenAI or hashing}
\STATE $\text{FAISS}_{\text{provider}} \gets \text{IndexFlatIP}(V_E)$
\RETURN $\{\text{RF}, \text{Cox}, \text{T-Learner}, \text{SHAP}, \text{FAISS}\}$
\end{algorithmic}
\end{algorithm}

\subsection{Resilient Multilingual RAG}
Each exit-interview narrative is tagged with a language code (EN /
ES / HI) and embedded with a provider chosen at request time. The
primary provider is OpenAI's
\texttt{text-embedding-3-small} (1{,}536-dim); the fallback is a
deterministic 256-dim MD5-bucketed hashing embedding. At inference,
the system performs a 250-millisecond reachability probe (DNS, TCP,
HTTPS HEAD on \texttt{api.openai.com}); on success the OpenAI index
is queried, on failure (firewall, DNS error, 5xx) the hashing index
is queried instead. Per-provider FAISS directories
(\texttt{faiss\_hr\_index\_multilingual/openai} versus
\texttt{faiss\_hr\_index\_multilingual/hashing}) prevent dimension
collisions. The dashboard surfaces a \texttt{provider} tag with
each result so the reviewer can see which path served the query.
This pattern --- cheap probe, per-provider stores, transparent
provider tag --- generalises beyond HR and is, we believe, a useful
deployment contribution in its own right.

\subsection{NL2SQL Agent}
We use the LangChain SQL agent pattern. The agent is initialised
with the SQLite schema, a read-only allow-list of four tables
(\texttt{employees}, \texttt{engagement\_pulse},
\texttt{exit\_interviews}, \texttt{workforce\_predictions}), a
few-shot example bank, and a self-correction loop: if the generated
query raises a SQL error, the error is fed back into the prompt for
one retry. We evaluate against a hand-curated 50-question gold set
spanning aggregate queries, joins, and filtered selections. When
the agent fails twice or the LLM is unreachable, a TF-IDF nearest-
question fallback returns the closest gold-set question and its
canonical SQL.

\subsection{Implementation and Reproducibility}
The system is implemented in Python~3.11 using
scikit-learn~\cite{pedregosa2011}, imbalanced-learn,
lifelines~\cite{davidson2019}, EconML~\cite{econml2019},
LangChain~\cite{langchain2023}, FAISS~\cite{johnson2019}, and
Streamlit. The full source, prompts, and evaluation harness are
released under MIT license. The codebase carries 31 \texttt{pytest}
tests across 8 test modules and a GitHub Actions workflow that runs
\texttt{ruff} and the test suite on every push. The codebase was
developed collaboratively with LLM coding assistants (GitHub
Copilot, Claude, GPT-4); the use of AI assistants is explicitly
documented in the project repository in line with emerging norms
for reproducible AI-assisted research, and the
\texttt{session\_logs/} directory captures verbatim prompts and
responses for traceability.

%==============================================================
\section{Evaluation}
%==============================================================

\subsection{Predictive Performance}
On a stratified 80/20 train/test split (Table~\ref{tab:perf}):

\begin{table}[h]
\centering
\caption{Predictive Performance on Held-Out Test Set}
\label{tab:perf}
\begin{tabular}{lcccc}
\toprule
\textbf{Bucket} & \textbf{Precision} & \textbf{Recall} & \textbf{F1} & \textbf{Support} \\
\midrule
Low      & 0.96 & 0.98 & 0.97 & 870 \\
Medium   & 0.71 & 0.62 & 0.66 & 41 \\
High     & 0.83 & 0.78 & 0.80 & 89 \\
\midrule
Macro    & 0.83 & 0.79 & 0.81 & 1{,}000 \\
\bottomrule
\end{tabular}
\end{table}

ROC-AUC is 0.91, PR-AUC is 0.74, and the calibration Brier score is
0.061. The PR-AUC is reported alongside ROC-AUC because PR is more
sensitive to the imbalanced positive class, and the Brier score
because deployment decisions condition on calibrated probabilities,
not just rank order.

\subsection{Survival-Model Goodness-of-Fit}
The Cox PH model attains a concordance index of 0.78 on the test
fold, with the proportional-hazards assumption holding for three of
the four covariates (\texttt{Tenure\_Months},
\texttt{Base\_Salary}, \texttt{Manager\_Tenure\_Months}). The
\texttt{Department\_Code} covariate shows a mild proportionality
violation, which we surface as a future-work item rather than
patch over.

\subsection{Uplift-Estimator Sanity Checks}
On the synthetic corpus where the data-generating process is known,
the T-Learner correctly recovers the sign of every CATE and
attains mean absolute error of 0.04 against the ground-truth treatment
effect. Table~\ref{tab:cate} reports the mean estimated CATE per
treatment arm on the high-risk cohort. We do not claim this transfers
to real observational data, which is precisely why the dashboard
surfaces the assumption.

\begin{table}[h]
\centering
\caption{Mean Estimated CATE per Treatment Arm (High-Risk Cohort)}
\label{tab:cate}
\begin{tabular}{lccc}
\toprule
\textbf{Arm} & \textbf{Mean $\hat{\tau}$} & \textbf{P5--P95} & \textbf{$n$ recommended} \\
\midrule
Compensation       & $-0.18$ & $[-0.31, -0.07]$ & 412 \\
Career-growth      & $-0.11$ & $[-0.22, -0.02]$ & 198 \\
Manager change     & $-0.07$ & $[-0.19, +0.02]$ & 87  \\
\bottomrule
\multicolumn{4}{l}{\footnotesize Negative $\hat{\tau}$ = expected risk reduction.}
\end{tabular}
\end{table}

\subsection{NL2SQL Agent Accuracy}
On the 50-question gold set, the agent achieves 86\% exact-match
SQL execution accuracy and 92\% semantic-match (correct answer,
syntactically different SQL). With the LLM disabled to force the
TF-IDF fallback, the system still returns a non-empty, ranked answer
on 100\% of queries, although the semantic-match accuracy drops to
58\%.

\subsection{Resilient-RAG Failover}
We benchmark the RAG path under three network conditions:
(i)~OpenAI reachable, (ii)~OpenAI unreachable (simulated firewall,
TCP RST), (iii)~OpenAI returning 5xx. In all three cases the
application remains responsive: the median end-to-end latency is
0.41~s, 0.34~s, and 0.55~s respectively. The provider tag in the
result row correctly reflects the active path on 100\% of queries.

\subsection{End-to-End Latency}
Median latency for a complete user query (NL2SQL $\rightarrow$ DB
execution $\rightarrow$ LLM-formatted answer) is 11.7~seconds versus
a baseline of 2.3~days for analyst-mediated queries estimated from
a small in-program survey, a reduction of approximately four orders
of magnitude.

\subsection{Manager Trust Study (Pilot)}
A pilot Likert-scale survey ($n{=}12$, MSBA cohort) compared
manager comfort acting on a black-box risk score versus a
RAG-grounded explanation. Mean trust rose from 2.4/5 to 4.1/5
($p < 0.01$, Wilcoxon signed-rank). We position this as
preliminary; a full IRB-approved study with HR practitioners is
proposed as future work.

\subsection{Ablation}
Table~\ref{tab:ablation} reports an ablation on three
deployment-relevant decisions. Removing the per-fold SMOTE
discipline (i.e.\ applying SMOTE before the split) inflates ROC-AUC
to 0.95 and PR-AUC to 0.81 --- a result we believe is widespread in
the published HR literature and which we consider a leakage
artefact. Removing the four-fifths gate has no effect on the model
metrics but makes the system non-compliant with NYC LL144. Removing
the provider fallback reduces firewall-condition availability from
100\% to 0\%.

\begin{table}[h]
\centering
\caption{Ablation of Deployment-Critical Decisions}
\label{tab:ablation}
\begin{tabular}{lccc}
\toprule
\textbf{Variant} & \textbf{ROC-AUC} & \textbf{PR-AUC} & \textbf{Avail.} \\
\midrule
Full system                    & 0.91 & 0.74 & 100\% \\
SMOTE before split             & 0.95* & 0.81* & 100\% \\
No four-fifths gate            & 0.91 & 0.74 & 100\% \\
No provider fallback           & 0.91 & 0.74 & 0\%   \\
\bottomrule
\multicolumn{4}{l}{\footnotesize *inflated by leakage; not a reliable estimate.}
\end{tabular}
\end{table}

%==============================================================
\section{Ethics, Fairness, and Limitations}
%==============================================================

\subsection{Fairness Audit}
The dashboard runs an EEOC four-fifths-rule audit on every load,
computing the disparate-impact ratio
$r = \min_g s_g \big/ \max_g s_g$, where $s_g$ is the high-risk
selection rate in group $g$. On the synthetic data, the audit
returns $r = 0.94$ across departments (within the 0.80
threshold). Table~\ref{tab:fair} reports the per-group selection
rates, sample sizes, and the resulting ratio. A failing audit pauses
downstream alerts for the affected group rather than auto-mitigating:
mitigation is a human decision, and the system's job is to surface
the disparity, log the audit row, and stop.

\begin{table}[h]
\centering
\caption{Four-Fifths Audit by Department (Synthetic)}
\label{tab:fair}
\begin{tabular}{lccc}
\toprule
\textbf{Department} & \textbf{$n$} & \textbf{Selection rate} & \textbf{Ratio vs.\ max} \\
\midrule
Engineering    & 1{,}640 & 0.21 & 1.00 \\
Manufacturing  & 1{,}210 & 0.20 & 0.95 \\
Sales          &   930 & 0.19 & 0.91 \\
Finance        &   620 & 0.20 & 0.96 \\
HR             &   320 & 0.20 & 0.94 \\
Legal          &   280 & 0.20 & 0.94 \\
\midrule
\multicolumn{3}{r}{\textbf{Disparate-impact ratio:}} & \textbf{0.94} \\
\bottomrule
\end{tabular}
\end{table}

\subsection{Identification Bias}
A subtle bias the model itself cannot fix is identification bias:
the human reviewer who sees \texttt{EMP\_2417} on a risk dashboard
and recognises Ravi from last quarter's all-hands anchors on memory
rather than evidence. The model is innocent --- the feature column
list is \texttt{[Tenure\_Months, Base\_Salary, Department\_Code,
Manager\_Tenure\_Months]} and never contains the identifier --- but
the user surface still enables the bias. AttriSense therefore
\emph{pseudonymizes at the dashboard boundary}: the UI shows a
deterministic salted-hash \texttt{Review\_ID} (format
\texttt{RV-NNNNNN}). The mapping from \texttt{Review\_ID} back to
\texttt{Emp\_ID} lives in a separate access-controlled table with
its own audit log; the salt rotates quarterly. We are not aware of
this pattern being explicitly documented in prior published HR
analytics systems and consider it a small but practically important
contribution.

\subsection{Synthetic-Data Caveat}
All performance numbers are reported on a synthetic corpus and may
not transfer to real enterprise data without re-tuning. The
synthetic generator was calibrated to public HR distributions but
cannot replicate firm-specific cultural dynamics, manager effects,
or industry seasonality. A real deployment must repeat the audit on
real data.

\subsection{Surveillance and Worker Dignity}
Predictive flight-risk systems carry real risk of being weaponized
against workers (preemptive demotion, exclusion from key projects).
We argue --- and the open-source license enforces --- that
AttriSense outputs should be visible to the employee in question on
request, and should never be used as the sole input to an adverse
employment action. The system's intended-use statement, published
in the repository, is explicit on this.

\subsection{Regulatory Alignment}
The system architecture is designed to satisfy NYC Local Law
144~\cite{nyclaw2023} bias-audit requirements and the EU AI Act's
high-risk-system documentation expectations: a model card, an
intended-use document, a fairness policy, and a session-log audit
trail are all part of the deliverable.

%==============================================================
\section{Lessons Learned in Production Hardening}
%==============================================================

The journey from a working notebook to a Streamlit dashboard a
hiring manager could deploy surfaced four lessons that we found
underrepresented in the literature and worth documenting:

\textbf{Per-fold SMOTE matters more than which oversampler.}
Across our experiments, the difference between SMOTE,
Borderline-SMOTE, and ADASYN was within 0.01 PR-AUC. The
difference between applying any of them \emph{before} versus
\emph{after} the train/test split was 0.07 PR-AUC --- and the
``before'' version is a silent leakage bug.

\textbf{Calibration is not optional.} A bucketed risk score
(High/Medium/Low) is only meaningful if the probability bins are
calibrated to observed frequencies. The system therefore reports
the Brier score and a calibration plot on every retraining.

\textbf{Identifiers leak even when features do not.} As
discussed in Section~VI, removing protected attributes from the
feature set is necessary but not sufficient. The reviewer's
short-term memory is a feature surface too.

\textbf{Provider fallback is cheap.} A 250-millisecond
reachability probe and a per-provider FAISS directory cost ten
lines of code and recover 100\% availability in firewalled
environments. We saw no published HR RAG system that did this.

%==============================================================
\section{Conclusion and Future Work}
%==============================================================
We have presented AttriSense, a unified open-source framework
that combines imbalanced-classification flight-risk prediction,
Cox proportional-hazards survival analysis, T-Learner causal
uplift, RAG-grounded qualitative explanation, and NL2SQL
data-lake access into a single workflow consumable by
non-technical HR users, and wraps the whole system in a
fairness-by-default discipline that includes a four-fifths-rule
audit gate and an identification-bias mitigation pattern. On a
5{,}000-employee synthetic corpus, the system achieves ROC-AUC
of 0.91, a calibration Brier of 0.061, a concordance index of
0.78, and reduces analyst-query latency by approximately four
orders of magnitude.

Future work includes: (i)~replacing the simulated treatment
assignments with observed treatments and adopting a DR-Learner
or double-machine-learning estimator; (ii)~intersectional
fairness audits (e.g.\ department $\times$ manager-tenure band);
(iii)~federated learning across multi-site enterprises;
(iv)~a snapshot-date column to support month-over-month drift
KPIs; and (v)~a full IRB-approved manager-trust field study at
scale with HR practitioners outside the academic setting.

\section*{Acknowledgments}
The author thanks the California State University, East Bay
Master of Science in Business Analytics program for the academic
foundation, and acknowledges the use of GitHub Copilot, Claude,
and GPT-4 as coding collaborators throughout the implementation.
All AI-assisted code was human-reviewed and is publicly
auditable in the project repository.

\section*{Code and Data Availability}
Source code, synthetic dataset generator, evaluation harness, and
the Streamlit dashboard are publicly available under MIT license:
\url{https://github.com/Dogiparthi-Sharada/AttriSense}.

%==============================================================
\begin{thebibliography}{00}

\bibitem{bersin2017} J.~Bersin, ``The employee experience platform:
A new category arrives,'' Deloitte Insights, 2017.

\bibitem{ibm2017} IBM Watson Analytics, ``HR Analytics Employee
Attrition \& Performance,'' IBM, 2017.

\bibitem{saradhi2011} V.~V.~Saradhi and G.~K.~Palshikar, ``Employee
churn prediction,'' \emph{Expert Systems with Applications},
vol.~38, no.~3, pp.~1999--2006, 2011.

\bibitem{sexton2005} R.~S.~Sexton, S.~McMurtrey,
J.~O.~Michalopoulos, and A.~M.~Smith, ``Employee turnover: A neural
network solution,'' \emph{Computers \& Operations Research},
vol.~32, no.~10, 2005.

\bibitem{yedida2018} R.~Yedida, R.~Reddy, R.~Vahi, R.~Jana,
A.~GV, and D.~Kulkarni, ``Employee attrition prediction,''
\emph{arXiv:1806.10480}, 2018.

\bibitem{chawla2002} N.~V.~Chawla, K.~W.~Bowyer, L.~O.~Hall, and
W.~P.~Kegelmeyer, ``SMOTE: Synthetic minority over-sampling
technique,'' \emph{Journal of Artificial Intelligence Research},
vol.~16, pp.~321--357, 2002.

\bibitem{chawla2003} N.~V.~Chawla, A.~Lazarevic, L.~O.~Hall, and
K.~W.~Bowyer, ``SMOTEBoost: Improving prediction of the minority
class in boosting,'' in \emph{PKDD}, 2003.

\bibitem{cox1972} D.~R.~Cox, ``Regression models and life-tables,''
\emph{Journal of the Royal Statistical Society B}, vol.~34, no.~2,
1972.

\bibitem{kalbfleisch2002} J.~D.~Kalbfleisch and R.~L.~Prentice,
\emph{The Statistical Analysis of Failure Time Data}, 2nd ed.,
Wiley, 2002.

\bibitem{kuhn2018} M.~Kuhn and K.~Johnson,
\emph{Applied Predictive Modeling}, Springer, 2018.

\bibitem{davidson2019} C.~Davidson-Pilon, ``lifelines: Survival
analysis in Python,'' \emph{Journal of Open Source Software},
vol.~4, no.~40, 2019.

\bibitem{wager2018} S.~Wager and S.~Athey, ``Estimation and
inference of heterogeneous treatment effects using random forests,''
\emph{Journal of the American Statistical Association}, vol.~113,
no.~523, 2018.

\bibitem{kennedy2020} E.~H.~Kennedy, ``Optimal doubly robust
estimation of heterogeneous causal effects,''
\emph{arXiv:2004.14497}, 2020.

\bibitem{kunzel2019} S.~R.~K\"unzel, J.~S.~Sekhon, P.~J.~Bickel, and
B.~Yu, ``Metalearners for estimating heterogeneous treatment
effects using machine learning,''
\emph{PNAS}, vol.~116, no.~10, 2019.

\bibitem{econml2019} Microsoft Research, ``EconML: A Python package
for ML-based heterogeneous treatment effects estimation,''
\url{https://github.com/microsoft/EconML}, 2019.

\bibitem{lewis2020} P.~Lewis et al., ``Retrieval-augmented
generation for knowledge-intensive NLP tasks,''
in \emph{NeurIPS}, 2020.

\bibitem{johnson2019} J.~Johnson, M.~Douze, and H.~J\'egou,
``Billion-scale similarity search with GPUs,''
\emph{IEEE Trans. Big Data}, vol.~7, no.~3, 2019.

\bibitem{androutsopoulos1995} I.~Androutsopoulos, G.~Ritchie, and
P.~Thanisch, ``Natural language interfaces to databases --- An
introduction,''
\emph{Natural Language Engineering}, vol.~1, no.~1, 1995.

\bibitem{rajkumar2022} N.~Rajkumar, R.~Li, and D.~Bahdanau,
``Evaluating the text-to-SQL capabilities of large language
models,'' \emph{arXiv:2204.00498}, 2022.

\bibitem{pourreza2023} M.~Pourreza and D.~Rafiei, ``DIN-SQL:
Decomposed in-context learning of text-to-SQL with self-correction,''
in \emph{NeurIPS}, 2023.

\bibitem{lundberg2017} S.~M.~Lundberg and S.-I.~Lee, ``A unified
approach to interpreting model predictions,'' in \emph{NeurIPS},
2017.

\bibitem{nyclaw2023} New York City Local Law 144 of 2021,
``Automated employment decision tools,'' effective July~2023.

\bibitem{pedregosa2011} F.~Pedregosa et al., ``Scikit-learn: Machine
learning in Python,'' \emph{JMLR}, vol.~12, pp.~2825--2830, 2011.

\bibitem{brown2020} T.~Brown et al., ``Language models are few-shot
learners,'' in \emph{NeurIPS}, 2020.

\bibitem{langchain2023} H.~Chase, ``LangChain,'' GitHub, 2023.
\url{https://github.com/langchain-ai/langchain}

\end{thebibliography}

\end{document}
